\section{系统架构与边界}

\begin{frame}{仓库架构：采集、构图、编排三者显式解耦}
\centering
\begin{tikzpicture}[
  box/.style={draw,rounded corners=2pt,minimum height=1.0cm,align=center,font=\small,thick},
  arr/.style={-{Latex[length=2mm]},thick,Navy}
]
\node[box,fill=SoftBlue,minimum width=3.1cm] (common) at (0,1.8) {\textbf{backend/common}\\主题常量与纯函数};
\node[box,fill=SoftGreen,minimum width=3.6cm] (intel) at (-4.25,0) {\textbf{backend/intel\_agent}\\采集、判断、覆盖、持久化};
\node[box,fill=Paper,minimum width=3.7cm] (console) at (0,0) {\textbf{backend/console}\\UI 编排、TraceSink、Database};
\node[box,fill=SoftAmber,minimum width=3.6cm] (kg) at (4.25,0) {\textbf{backend/ctinexus\_kg}\\消费 item，生成 item-KG};
\draw[arr] (common) -- (intel);
\draw[arr] (common) -- (kg);
\draw[arr] (console) -- (intel);
\draw[arr] (console) -- (kg);
\draw[arr,dashed,bend right=34] (intel.south) to node[below,font=\scriptsize,align=center]{run\_id + item DTO / 持久化产物} (kg.south);
\end{tikzpicture}

\vspace{0.25em}
\begin{alertblock}{架构边界}
\centering
\small \pathcode{intel_agent} 与 \pathcode{ctinexus_kg} 禁止互相 import；KG 失败不得阻断采集主流程。
\end{alertblock}

\smallsource{\pathcode{AGENTS.md}；\pathcode{backend/common/topics.py:1-5}}
\end{frame}

\begin{frame}{部署架构：一个在线控制台 + 一个批处理采集器}
\begin{columns}[T,onlytextwidth]
\column{0.62\textwidth}
\centering
\resizebox{\linewidth}{!}{%
\begin{tikzpicture}[
  n/.style={draw,rounded corners,align=center,minimum height=0.8cm,font=\scriptsize,thick},
  a/.style={-{Latex[length=2mm]},thick,Navy}
]
\node[n,fill=SoftBlue,minimum width=2.4cm] (user) {浏览器 / CLI};
\node[n,fill=SoftGreen,minimum width=3.2cm,right=0.8cm of user] (web) {web container\\Gradio + 统一镜像};
\node[n,fill=SoftAmber,minimum width=3.0cm,below=0.75cm of web] (mongo) {MongoDB\\container:27017\\host:27018};
\node[n,fill=Paper,minimum width=3.2cm,above=0.75cm of web] (batch) {intel-agent profile\\full / incremental};
\node[n,fill=SoftRed,minimum width=2.8cm,right=0.8cm of web] (api) {Claude/GLM/DeepSeek\\Anthropic-compatible};
\draw[a] (user) -- (web);
\draw[a] (user.north east) -- (batch.west);
\draw[a,<->] (web) -- (mongo);
\draw[a,<->] (batch) -- (mongo);
\draw[a,<->] (web) -- (api);
\draw[a,<->] (batch) -- (api);
\end{tikzpicture}
}
\column{0.35\textwidth}
\begin{block}{容器安全前提}
镜像以 root 运行时必须设置 \pathcode{IS_SANDBOX=1}，否则 SDK 的 \pathcode{bypassPermissions} 会被拒绝并触发降级。
\end{block}
\begin{block}{持久层}
优先 Mongo；未启用时写 \pathcode{data/intel_runs/*.json}。trace 始终可落 JSONL。
\end{block}
\end{columns}
\smallsource{\pathcode{docker-compose.yml}；\pathcode{Dockerfile.intel_agent}；根目录项目指南}
\end{frame}

\begin{frame}{控制面与数据面：LLM 提议，Python 保留权威}
\centering
\begin{tikzpicture}[
  lane/.style={draw,rounded corners,minimum width=13.8cm,minimum height=1.05cm,align=left,inner sep=8pt},
  chip/.style={draw,rounded corners=2pt,minimum height=0.52cm,align=center,font=\scriptsize,inner sep=4pt},
  a/.style={-{Latex[length=1.8mm]},thick,Navy}
]
\node[lane,fill=SoftBlue] (control) at (0,1.8) {};
\node[anchor=west,font=\small\bfseries,text=Navy] at (-6.65,2.12) {控制面};
\node[chip,fill=white] (ctx) at (-4.5,1.68) {状态摘要};
\node[chip,fill=white,right=0.38cm of ctx] (llm) {LLM 生成候选};
\node[chip,fill=white,right=0.38cm of llm] (ucb) {确定性评分/UCB};
\node[chip,fill=white,right=0.38cm of ucb] (guard) {schema + hooks};
\node[chip,fill=white,right=0.38cm of guard] (stop) {critic/停止};
\draw[a] (ctx)--(llm);\draw[a] (llm)--(ucb);\draw[a] (ucb)--(guard);\draw[a] (guard)--(stop);

\node[lane,fill=SoftGreen] (data) at (0,0.2) {};
\node[anchor=west,font=\small\bfseries,text=Green] at (-6.65,0.52) {数据面};
\node[chip,fill=white] (source) at (-4.85,0.08) {NVD/arXiv/OSV/CISA};
\node[chip,fill=white,right=0.38cm of source] (norm) {统一 RawIntelItem};
\node[chip,fill=white,right=0.38cm of norm] (rel) {相关性/主题};
\node[chip,fill=white,right=0.38cm of rel] (dedup) {事件去重/coverage};
\node[chip,fill=white,right=0.38cm of dedup] (store) {Mongo/JSON/trace};
\draw[a] (source)--(norm);\draw[a] (norm)--(rel);\draw[a] (rel)--(dedup);\draw[a] (dedup)--(store);
\draw[a,dashed] (store.north) -- node[right,font=\tiny]{blackboard} (ctx.south);
\draw[a,dashed] (guard.south) -- node[right,font=\tiny]{approved call} (source.north);
\end{tikzpicture}

\vspace{0.4em}
\begin{block}{关键原则}
模型可以提出 query、source、operator 与继续/停止理由；预算、候选合法性、权威时间窗、去重、持久化和硬安全阀不交给模型。
\end{block}
\end{frame}

\begin{frame}{端到端数据流：从候选到可追溯 item-KG}
\centering
\begin{tikzpicture}[
  flowstep/.style={draw,rounded corners=2pt,minimum width=2.0cm,minimum height=0.76cm,align=center,font=\scriptsize,thick},
  a/.style={-{Latex[length=1.7mm]},thick,Navy},
  node distance=0.42cm
]
\node[flowstep,fill=SoftBlue] (cand) {QueryCandidate\\严格 schema};
\node[flowstep,fill=SoftBlue,right=of cand] (approve) {UCB portfolio\\+ hooks};
\node[flowstep,fill=SoftGreen,right=of approve] (call) {源 API\\分页/重试};
\node[flowstep,fill=SoftGreen,right=of call] (raw) {RawIntelItem\\统一字段};
\node[flowstep,fill=SoftAmber,right=of raw] (rel) {relevance\\topic scores};
\node[flowstep,fill=SoftAmber,below=0.65cm of raw] (gap) {CoverageGap /\\CorpusGap};
\node[flowstep,fill=Paper,left=of gap] (bb) {Blackboard\\outcome/checkpoint};
\node[flowstep,fill=Paper,left=of bb] (persist) {Mongo/JSON\\trace/playbook};
\node[flowstep,fill=SoftRed,right=of gap] (kg) {item-KG\\独立下游};
\draw[a] (cand)--(approve);\draw[a] (approve)--(call);\draw[a] (call)--(raw);\draw[a] (raw)--(rel);
\draw[a] (rel)--(gap);\draw[a] (gap)--(bb);\draw[a] (bb)--(persist);\draw[a,dashed] (persist)--(kg);
\draw[a,dashed,bend left=25] (bb.north) to node[above,font=\tiny]{下一轮状态} (cand.south);
\end{tikzpicture}

\vspace{0.7em}
\begin{itemize}
  \item 自然键 \pathcode{item_id} 在 run 内和全局 Mongo 两层去重，\pathcode{run_ids} 保存溯源。
  \item item-KG 只消费 item 的 \pathcode{raw_text}（空则 summary）；\pathcode{source_uri} 仅作 provenance。
\end{itemize}
\end{frame}

\begin{frame}{Blackboard：闭环控制的单一运行态}
\scriptsize
\begin{tabularx}{\textwidth}{p{0.22\textwidth} X p{0.29\textwidth}}
\toprule
状态组 & 关键字段 & 控制意义 \\
\midrule
任务与策略 & run\_goal, run\_mode, collection\_strategy & 区分 full legacy 与 incremental adaptive \\
预算 & max\_rounds, max\_seconds, max\_api\_calls & hooks 与 controller 的硬边界 \\
观测 & raw\_items, query\_history, query\_outcomes & novelty/noise/duplicate/reward 的事实来源 \\
覆盖 & coverage\_gaps, corpus\_gaps, run\_gaps & legacy 配额与有效证据是两套不同口径 \\
连续性 & source\_checkpoints, sdk\_session\_id & 增量水位、分页恢复与 SDK resume \\
学习 & candidate\_topics, extended\_trends, reflection\_notes & 新主题、burst、critic 解释 \\
审计 & action\_history, errors, metrics & 动作、失败、调用数与时间 \\
\bottomrule
\end{tabularx}

\smallsource{\pathcode{backend/intel_agent/schemas.py:196-327}}
\end{frame}

\begin{frame}{主题本体：18 Core 驱动覆盖，15 Extended 负责发现}
\begin{columns}[T,onlytextwidth]
\column{0.54\textwidth}
\begin{block}{Core（18）}
\scriptsize
prompt injection；jailbreak；agent tool abuse；data leakage；model supply chain；RAG poisoning；insecure output handling；excessive agency；model denial of service；unbounded consumption；model backdoor；agent loop/cascading failure；multi-agent communication attack；unsafe code generation/execution；sandbox escape；plugin/MCP trust；authentication/authorization bypass；unsafe file handling。
\end{block}
\column{0.43\textwidth}
\begin{block}{Extended（15）}
\scriptsize
训练/微调投毒、对抗样本、模型抽取、成员推断、模型反演、训练数据抽取、system prompt 泄漏、context isolation、memory poisoning、agent impersonation、SSRF、内部网侦察、虚假信息/引用、日志监控不足、模型路由/降级。
\end{block}
\end{columns}

\vspace{0.6em}
\begin{alertblock}{不能把 taxonomy 当作完备本体}
它是可搜索标签集合，不是 OWASP/MITRE 的一一映射；同一 item 可多标签，adaptive 用 topic credit 归一化避免重复全额计数。
\end{alertblock}
\smallsource{\pathcode{backend/common/topics.py:11-55}；外部校准参照 \cite{owaspllm2025,mitreatlas}}
\end{frame}
